\documentclass[10pt,a4paper]{amsart}
\usepackage[margin=19mm]{geometry}
\usepackage{amsmath,amssymb,mathtools}
\usepackage[scr=rsfs,cal=euler]{mathalfa}
\usepackage{xfrac}
\usepackage[unicode,hidelinks,bookmarks=false]{hyperref}
\newcommand{\ie}{i.e.\ }
\newcommand{\cf}{cf.\ }
\newcommand{\resp}{respectively\ }
\newcommand{\Subsection}{\subsection}
\providecommand{\nobreakdash}{\nobreak\hbox{-}}
\setlength{\emergencystretch}{2em}
\multlinegap0pt
\usepackage{bm}
\usepackage{bbm}
\usepackage{dsfont}
\usepackage{xspace}
\usepackage{cancel}
\usepackage{mathtools}
\usepackage{amsthm}
\makeatletter
\@ifundefined{thm}{\newtheorem{thm}{Thm}}{}
\@ifundefined{proposition}{\newtheorem{proposition}[thm]{Proposition}}{}
\makeatother
\newcommand {\T}  { {\mathbb{T}} }
\title[A cross-diffusion system with independent drifts and fast di]{A cross-diffusion system with independent drifts and fast diffusion}
\author{Charles Elbar, Filippo Santambrogio}
\date{Source: arXiv:2510.07937v1 (CC BY 4.0)}
\begin{document}
\maketitle

\noindent\emph{Source and licence.} A cross-diffusion system with independent drifts and fast diffusion. Version arXiv:2510.07937v1 (CC BY 4.0), distributed under the Creative Commons Attribution 4.0 International licence, as recorded in the arXiv licence metadata for this exact version and shown on the abstract page https://arxiv.org/abs/2510.07937v1. Full rights notice: licenses/arxiv-2510.07937-CC-BY-4.0.txt.\par\smallskip

\noindent\textit{Editorial scope.} Counted unit: the Proposition whose source label is \texttt{prop:spatial\_S} in the article (source0.tex lines 424--431, source label \texttt{prop:spatial\_S}), delivered together with the article's own complete original proof of it (source0.tex lines 433--471, one \texttt{proof} environment of 39 lines). No mathematics is written, repaired or re-derived: statement and proof are the article's own text line for line. Required context retained uncounted: eq:prelim\_dS (lines 415--420). Every definition, display and label of those lines is retained unchanged. Omitted from this package and not redistributed: 1--414, 421--423, 432--432, 472--855. Nothing inside a retained excerpt is omitted or altered: every retained body is delivered exactly as the article's lines are written, so no omission receipt of any kind accompanies this package. Citations: the retained excerpts carry no citation command at all, so no bibliography entry of the article is transcribed and no reference is redistributed. Reference adaptation: none was needed and none was made; every reference inside the delivered unit resolves inside it, so no unresolved reference or citation is produced. Printed slips kept literal: no printed word, symbol, hypothesis or number of the counted statement or of its proof was altered. Layout and class: the article's own class is \texttt{article} with packages including inputenc, fontenc, geometry, amsmath, amssymb, mathtools, amsthm, mathrsfs, graphicx, xcolor, babel, bbm, appendix, caption; the wrapper is the lane's pinned \texttt{amsart} editorial wrapper at 10pt on A4, which restates the theorem-like environments the retained text needs and adds the article's own macro definitions that the retained text needs (\texttt{\textbackslash T}), with extra wrapper packages (bm, bbm, dsfont, xspace, cancel, mathtools). The delivered numbering is the unit's own. Assets: no retained span contains a figure environment, a graphics-inclusion command, a PDF-page inclusion, a table or a TikZ picture; no image file is copied into this package, the package declares no asset dependency and no image file is redistributed with it. Licence: version arXiv:2510.07937v1 of this article is distributed under the Creative Commons Attribution 4.0 International licence, as recorded in the arXiv licence metadata for this exact version and shown on the abstract page https://arxiv.org/abs/2510.07937v1. A full rights notice is delivered at licenses/arxiv-2510.07937-CC-BY-4.0.txt. Characters: every retained excerpt is pure ASCII, so the delivered engine, XeLaTeX with its default Latin Modern text font, reports no missing character.
\begin{equation}\label{eq:prelim_dS}
\partial_t S 
 =  
\partial_{xx} \Phi(S)
 + \partial_x \left(\rho\partial_x V + \mu\partial_x W \right).
\end{equation}

\begin{proposition}[Spatial estimate on $S$]\label{prop:spatial_S}
We have the following estimates on $S$:
\begin{align*}
&S\in L^{\infty}((0,T); L^{1}(\T)), \quad \partial_x S^{\alpha-\frac{1}{2}}\in L^{2}((0,T)\times \T),\quad S^{\alpha-\frac{1}{2}}\in L^{2}((0,T); L^{\infty}(\T)),\\
&\partial_x \log S\in L^{2}((0,T)\times \T), \quad \partial_x S^{\frac{\alpha}{2}}\in L^{2}((0,T)\times \T) \quad S^{\frac{\alpha}{2}} \in L^{2}((0,T); L^{\infty}(\T)),\\
&S\in L^{2-\alpha}((0,T)\times \T),  \quad \partial_x S^{1-\alpha}\in L^{1}((0,T)\times \T), \quad  S^{1-\alpha}\in L^{1}((0,T); L^{\infty}(\T)). 
\end{align*}
\end{proposition}

\begin{proof}
We first recall that the nonnegativity of both $\rho$ and $\mu$ can be obtained from the transport equation structure and the nonnegativity of the initial conditions. The first estimate is obtained after integrating~\eqref{eq:prelim_dS} in space and using periodic boundary conditions as well as the integrability of the initial conditions. 
We recall that $\Phi(S)= S^{\alpha}$ for $0< \alpha\le 1$. We only do the proof of the second estimate for $\alpha<1$, the proof for $\alpha=1$ being similar by considering the dissipation of $\frac{d}{dt}\int_{\T}S\log S$.
 Using~\eqref{eq:prelim_dS}, $\rho,\mu\le S$,  $|\partial_x V|, |\partial_x W|\le C$ and Young's inequality we obtain for $0\le \beta\le  1$

\begin{multline*}
\frac{d}{dt}\int_{\T}-S^{\beta}  + \alpha\beta(1-\beta)\int_{\T}S^{-3+\alpha+\beta}(\partial_x S)^2 \\= \beta(\beta-1)\int_{\T}\rho\partial_x V S^{\beta-2}\partial_xS+ \beta(\beta-1)\int_{\T}\mu\partial_x W S^{\beta-2}\partial_xS\\
\le \frac{\alpha\beta(1-\beta)}{2}\int_{\T}S^{-3+\alpha+\beta}(\partial_x S)^2 + C \int_{\T}S^{\beta+1-\alpha}.
\end{multline*}

Therefore we have

\begin{equation*}
\frac{d}{dt}\int_{\T}-S^{\beta} + \frac{2\alpha\beta(1-\beta)}{(\alpha+\beta-1)^2}\int_{\T}\left(\partial_x S^{\frac{\alpha+\beta-1}{2}}\right)^2 \le  C \int_{\T}S^{\beta+1-\alpha}.
\end{equation*}
Note that in this expression, the term in the middle reads $\frac{\alpha\beta(1-\beta)}{2}\int_\T \left(\partial_x \log S\right)^2$ when $\alpha+\beta=1$.

As $S\in L^{\infty}((0,T); L^{1}(\T))$, to obtain a priori estimates from this dissipation we first choose $\beta=\alpha$  and we obtain the second estimate after integrating in time and using also the integrability of the initial condition. The third estimate is a consequence of the first two estimates. Choosing then $\beta=1-\alpha$ and using the seventh estimate that is proved independently below we obtain the fourth estimate (indeed, for the validity of the fourth estimate we would need $S\in L^{2-2\alpha}$ and the seventh provides $S\in L^{2-\alpha}$; of course, when $\alpha>1/2$ this argument is not necessary).
To obtain the fifth estimate we compute the sum of the dissipation of each entropies: 
\begin{equation*}
\frac{d}{dt}\int_{\T}\rho\log\rho +\mu\log\mu + \int_{\T}f''(S)|\partial_xS|^2 = \int_{\T} \rho \partial_{xx}V +\int_{\T}\mu\partial_{xx}W\le C.   
\end{equation*}
The last estimate of this inequality is a consequence of mass conservation and the assumptions on $V,W$. Since $f''(S)|\partial_x S|^2 = C|\partial_x S^{\frac{\alpha}{2}}|^2$ we obtain the fifth estimate. We deduce the sixth estimate from the fifth.
\iffalse
To prove the seventh and eighth estimate we take $\varphi$ such that $-\partial_{xx}\varphi = S-\bar{S}$ where $\bar{S}$ is the mass of $S$. We obtain 
$$
\frac{d}{dt}\int_{\T}|\partial_ x\varphi|^2 + \int_{\T}S^{1+\alpha} = -\int_{\T}S Q\partial_x\varphi
$$
where $Q= \frac{\rho\partial_x V + \mu \partial_x W}{S}$ is bounded in $L^{\infty}((0,T)\times\T)$. Moreover $|\partial_x\varphi|$ is bounded in $L^{\infty}((0,T);L^{1}(\T))$ by definition of $\varphi$ and $S\in L^{\infty}((0,T); L^{1}(\T))$. Therefore Gronwall's lemma yields the seventh eighth estimate.
\fi
From the eighth estimate we obtain the ninth, so it only remains to prove the seventh and eighth estimates. By using the previous estimates and Jensen's inequality we obtain:

\begin{align*}
\int_{0}^{T}\int_{\T}S^{2-\alpha}&= \int_{0}^{T}\int_{\T}SS^{1-\alpha}
\lesssim  \int_{0}^{T}\|S^{1-\alpha}\|_{L^{\infty}} \lesssim \int_{0}^{T}\int_{\T}\left|\partial_x S^{1-\alpha}\right|\lesssim \int_{0}^{T}\int_{\T}S^{1-\frac{3\alpha}{2}}|\partial_x S^{\frac{\alpha}{2}}|\\
&\lesssim \sqrt{\int_{0}^{T}\int_{\T}S^{2-3\alpha}}\sqrt{\int_{0}^{T}\int_{\T}|\partial_x S^{\frac{\alpha}{2}}|^2} \lesssim \left(\int_{0}^{T}\int_{\T}S^{2-\alpha}\right)^{\frac{1}{2}\frac{2-3\alpha}{2-\alpha}}
\end{align*}
From these inequalities we obtain the last two estimates: the seventh because $\frac{1}{2}\frac{2-3\alpha}{2-\alpha}<1$ and then the eighth one using the seventh estimate. 
\end{proof}

\end{document}
